\begin{textsample}{POS Dim 1 – vem\_ed\_09 – Score 14.00 – vem\_ed\_09/t038.txt}  \label{ex:f1_pos_085}
BOAS PRÁTICAS Competências culturais: além da \textbf{disciplina} Mestre em Economia Aplicada pela Esalq / USP, onde se graduou em Engenharia Agronômica, Ricardo Bertoni Pompeu \textbf{é} professor da Fatec Americana desde 2010 e atua com COIL / PCIs há quatro anos.
Pompeu \textbf{compartilha} um pouco de sua experiência no relato a \textbf{seguir}: Em 2017, participei de um megaprojeto com Eva Haug e Ariane Hoekstra, da Amsterdam University of Applied Sciences (AUAS / Holanda) e Osvaldo Succi Junior (CPS), Cristine Moraes (Fatec Piracicaba), Ana Teresa Colenci Trevelin (Fatec São Carlos), Carlos Amaral Moreira e eu (Fatec Americana).
Durante 7 semanas, 70 alunos estudaram \textbf{estratégias} de comunicação de marcas de suco de frutas, cerveja, refrigerante, chocolate, bolacha e sorvete nos dois países.
Em 2018, nosso projeto envolveu os professores Arun Pillutla (St.
Ambrose University / EUA), Carlos Moreira, Osvaldo Succi Junior e eu.
O PCI abordou a \textbf{competência} de dar e receber feedback, com a Metodologia Ativa de Aprendizado Baseado em Problemas.
Os alunos receberam uma situação problema e textos teóricos sobre feedback, incluindo o aspecto cultural dessa \textbf{competência}.
Em 2019 tivemos um novo PCI com a Holanda, com Eva Haug e Brechtine Detmar (AUAS) e Osvaldo Succi Junior (CPS / Cesu), Carlos Moreira e eu.
Optamos por um modelo diferente, convidando estudantes de vários semestres de Gestão Empresarial.
Nesse projeto, os 57 alunos brasileiros e holandeses abordaram a mesma situação problema sobre feedback (do PCI de 2018).
Durou 6 semanas, com grande enfoque na \textbf{parte} cultural da comunicação, da liderança e da gestão de equipes.
Agora, estou desenvolvendo um projeto com o Symbiosis Centre for Management Studies (Índia).
As \textbf{expectativas} \textbf{são} as melhores \textbf{possíveis}, pois percebo uma química muito grande entre os professores Nehajoan Panackal (Symbiosis), Rafaeli Begalli Danilo Sorroce (Fatec Sumaré) e eu.
Não \textbf{importa} se um projeto durar 4 ou 7 semanas, \textbf{é} importante uma preparação estratégica prévia, para entender as \textbf{expectativas}, interesses e ansiedades dos professores.
Ao mesmo tempo, compreender os alunos: \textbf{quantos} \textbf{são}, \textbf{conhecimento} do \textbf{idioma}, idade, o \textbf{que} já estudaram, semestre de curso.
Outro \textbf{fator} importante \textbf{é} a afinidade entre as \textbf{disciplinas} dos professores envolvidos.
No meu caso, leciono Gestão de Pessoas e Comportamento Organizacional, \textbf{disciplinas} \textbf{que} possuem aderência com várias áreas da administração.
A principal dica \textbf{é} não se ater somente à ementa da \textbf{disciplina}, mas ter uma visão \textbf{ampla} sobre as \textbf{competências} a \textbf{serem} desenvolvidas: cultura, comunicação, gestão de \textbf{conflitos} e trabalho em equipe.
Daí a relevância do trabalho da equipe dos PCIs / Cesu, ao encontrar parcerias personalizadas para \textbf{cada} professor ou professora \textbf{que} se interessar em participar de um PCI.
Leia a íntegra do depoimento em: https://cesu.cps.sp.gov.br/professor-da-fatec-americana-compartilha-experiencia-em-pcis/

% matched lemmas: amplo, cada, compartilhar, competência, conflito, conhecimento, disciplina, estratégia, expectativa, fator, idioma, importar, parte, possível, quanto, que, seguir, ser
\end{textsample}
